% ============================================================
%  Capitolo 4 — Lo scheduler carbon-aware
% ============================================================

\chapter{Lo scheduler carbon-aware}
\label{cap:scheduler}

Il Capitolo~\ref{cap:ProgettazioneRealizzazione} ha descritto
l'architettura del sistema, il modello di funzione logica con varianti
e i meccanismi di aggiornamento adattivo dei profili energetici e di
esplorazione LCB. Il presente capitolo introduce il componente
decisionale che costituisce il nucleo dell'estensione proposta: lo
scheduler carbon-aware.

Lo scheduler traduce il segnale ambientale --- la carbon intensity della
zona geografica del nodo al momento dell'invocazione --- in una decisione
di selezione della variante da eseguire. La logica decisionale si
articola in tre fasi: la normalizzazione della carbon intensity corrente
in un peso $\lambda \in (0,1)$ mediante una sigmoide calibrata per zona
(\S\,\ref{sec:sigmoide}); la costruzione della frontiera di Pareto nello
spazio energia-errore (\S\,\ref{sec:pareto}); la selezione della
variante Pareto-ottimale tramite scalarizzazione lineare pesata da
$\lambda$ (\S\,\ref{sec:scalarizzazione}).


% ============================================================
\section{La carbon intensity come segnale decisionale}
\label{sec:sigmoide}
% ============================================================

La carbon intensity (CI) esprime la quantità di
CO\textsubscript{2}~equivalente emessa per ogni kilowattora prodotto
dalla rete elettrica locale, misurata in
g\,CO\textsubscript{2}eq/kWh. Questo valore riflette la composizione
del mix energetico di una zona geografica in un dato istante: una rete
alimentata prevalentemente da fonti rinnovabili presenta valori di CI
ridotti, mentre una rete dipendente da combustibili fossili mostra
valori elevati.

Nel contesto di questo lavoro, la CI viene utilizzata come segnale
per derivare il parametro di trade-off $\lambda \in (0,1)$ che governa
la selezione della variante (Sezione~\ref{sec:pareto}): quando
l'energia disponibile è a bassa intensità carbonica, il sistema può
privilegiare l'accuratezza del risultato senza un costo ambientale
significativo; quando la rete è carbon-intensive, è preferibile
ridurre i consumi a scapito dell'accuratezza.

Tuttavia, il valore assoluto della CI non è sufficiente per prendere
questa decisione: una CI di 100\,g\,CO\textsubscript{2}/kWh
rappresenta un picco anomalo per la Norvegia, ma è del tutto ordinaria
per la Spagna. Stabilire se l'energia è più o meno ``green'' in un
dato momento richiede quindi un riferimento contestuale, ovvero la
distribuzione storica della CI per la zona geografica considerata.

È necessario pertanto un modello di normalizzazione che trasformi la
CI corrente in un peso $\lambda \in (0,1)$ relativo al contesto storico
della zona. La semantica adottata è la seguente:
%
\begin{itemize}
    \item $\lambda \to 1$: la CI corrente è bassa rispetto alla
    distribuzione storica della zona (rete elettrica relativamente
    \emph{pulita}); il sistema privilegia l'accuratezza del risultato;
    \item $\lambda \to 0$: la CI corrente è elevata (rete
    \emph{sporca}); il sistema privilegia l'efficienza energetica.
\end{itemize}

% ------------------------------------------------------------
\subsection{La funzione sigmoide calibrata per zona}
\label{sec:sigmoide:calibrazione}
% ------------------------------------------------------------

Per tradurre la CI corrente in un valore di $\lambda$ relativo al
contesto storico della zona, si adotta una funzione sigmoide logistica
nella forma:
%
\begin{equation}
    \lambda(\mathrm{CI}) \;=\;
        \frac{1}
             {1 + \exp\!\left(k \cdot
                 \dfrac{\mathrm{CI} - \mu_z}{s_z}\right)}
    \label{eq:sigmoide}
\end{equation}
%
dove $\mu_z$ è la media storica della CI per la zona geografica~$z$
(punto di simmetria in cui $\lambda = 0{,}5$), $s_z$ è un fattore di
scala derivato dalla dispersione storica locale, e $k > 0$ è un
parametro di pendenza che controlla la velocità di transizione.
La scelta di questa famiglia di funzioni è motivata da tre proprietà
desiderabili nel contesto dello scheduling:
%
\begin{itemize}
    \item \textbf{Monotonicità decrescente}: $\lambda$ decresce
    monotonicamente al crescere di CI, associando in modo coerente
    condizioni di rete più carbon-intensive a un peso maggiore
    sull'efficienza energetica.

    \item \textbf{Saturazione}: $\lambda$ tende asintoticamente a~$1$
    per CI basse e a~$0$ per CI alte.

    \item \textbf{Sensibilità contestuale}: i parametri $\mu_z$ e $s_z$
    sono calibrati a partire dalla distribuzione storica della CI nella
    zona $z$, permettendo alla sigmoide di adattarsi al contesto locale.
    In questo modo, il valore di $\lambda$ esprime una misura relativa
    rispetto alla norma locale, rendendo confrontabili le decisioni di
    scheduling tra aree geografiche diverse, anche quando i livelli
    assoluti di carbon intensity differiscono significativamente.
\end{itemize}

\medskip
\noindent\textbf{Calibrazione dei parametri.}
I parametri $\mu_z$, $s_z$ e $k$ sono stimati a partire dai
dati storici di carbon intensity per ciascuna zona. La procedura
è la seguente:
%
\begin{enumerate}
    \item $\mu_z$ è posto pari alla \emph{media} storica della CI per
    la zona, garantendo $\lambda(\mu_z) = 0{,}5$, ovvero che il punto
    di bilanciamento neutro corrisponda al valore tipico della zona;

    \item $s_z$ è ricavato dal quinto e novantacinquesimo percentile
    della distribuzione storica:
    %
    \begin{equation}
        s_z \;=\;
            \frac{\mathrm{CI}_{95} - \mathrm{CI}_5}
                {2\,\ln\!\left(\frac{1-\alpha}{\alpha}\right)}
        \label{eq:s_calibrazione}
    \end{equation}
    %
    con $\alpha = 0{,}01$. L'utilizzo dei percentili estremi consente
    di catturare l'intervallo di variabilità tipica della CI nella zona,
    trascurando i valori anomali. In particolare, il 90\,\% centrale
    della distribuzione storica viene mappato nell'intervallo
    $[\alpha,\, 1-\alpha]$ di $\lambda$, garantendo che la funzione
    sia sufficientemente sensibile alle variazioni reali della CI,
    senza risultare distorta dalla presenza di outlier.

    \item $k$ è un parametro di pendenza configurabile per zona:
    valori più alti producono una transizione più netta attorno a
    $\mu_z$, valori più bassi una risposta più graduale.
\end{enumerate}

\medskip
\noindent\textbf{Calcolo a runtime.}

\begin{algorithm}[H]
\caption{Calcolo del peso $\lambda$ dalla carbon intensity corrente}
\label{alg:lambda}
\begin{algorithmic}[1]
\Require zona geografica $z$, soglia di fallback $\lambda_0 = 0{,}5$
\Ensure $\lambda \in (0, 1)$

\State $\mathrm{CI} \leftarrow \textsc{GetCarbonIntensity}(z)$
    \Comment{API ElectricityMaps}
\If{$\mathrm{CI}$ non disponibile}
    \State \Return $\lambda_0$
\EndIf

\State $\mu_z \leftarrow \textsc{GetMean}(z)$
\State $s_z \leftarrow \dfrac{\mathrm{CI}_{95}(z) - \mathrm{CI}_5(z)}
                              {2\,\ln\!\left(\dfrac{1-\alpha}{\alpha}\right)}$
\State $k \leftarrow \textsc{GetSlopeParam}(z)$

\State \Return $\dfrac{1}{1 + \exp\!\left(k \cdot \dfrac{\mathrm{CI} - \mu_z}{s_z}\right)}$
\end{algorithmic}
\end{algorithm}


Al momento dell'invocazione, lo scheduler recupera la CI
corrente tramite l'API di ElectricityMaps~\cite{electricitymaps2024}
, legge i parametri della zona, e calcola $\lambda$ secondo
l'Equazione~\eqref{eq:sigmoide}, come descritto nell'algoritmo \autoref{alg:lambda}. In caso di indisponibilità dell'API,
il sistema degrada a $\lambda = 0{,}5$ (bilanciamento neutro tra
accuratezza ed efficienza energetica).

% ============================================================
\section{La frontiera di Pareto nello spazio energia-errore}
\label{sec:pareto}
% ============================================================

L'insieme delle varianti registrate per una funzione logica definisce
uno spazio di soluzioni bidimensionale caratterizzato dal consumo
energetico stimato $E(v)$ e dal punteggio di errore
$\mathit{Err}(v)$, descritto nella
Sezione~\ref{sec:modello-varianti}. Poiché tra queste due grandezze
esiste un trade-off strutturale, la selezione ottimale non può
minimizzarne simultaneamente entrambe.

Il sistema restringe quindi la selezione alle sole varianti
\emph{Pareto-ottimali}: quelle per cui non esiste alcuna alternativa
migliore in entrambe le dimensioni.

\medskip
\noindent\textbf{Definizione.}
Sia $V = \{v_1, \ldots, v_n\}$ l'insieme delle varianti disponibili.
Una variante $v_i$ è \emph{dominata} da $v_j$ se:
%
\begin{equation}
    E(v_j) \leq E(v_i) \;\wedge\;
    \mathit{Err}(v_j) \leq \mathit{Err}(v_i)
    \;\wedge\;
    \bigl(E(v_j) < E(v_i) \;\vee\;
          \mathit{Err}(v_j) < \mathit{Err}(v_i)\bigr)
    \label{eq:dominanza}
\end{equation}

La frontiera di Pareto $\mathcal{P} \subseteq V$ è l'insieme delle
varianti non dominate:
%
\begin{equation}
    \mathcal{P} = \bigl\{\, v \in V \;\big|\;
        \nexists\, v' \in V : v' \text{ domina } v \,\bigr\}
    \label{eq:pareto}
\end{equation}

Geometricamente, $\mathcal{P}$ costituisce il fronte
inferiore-sinistro della nuvola di punti nello spazio
$(E,\,\mathit{Err})$: ogni variante sulla frontiera rappresenta un
compromesso efficiente, nel senso che ridurre una delle due grandezze
implica necessariamente aumentare l'altra.

\medskip
\noindent\textbf{Costruzione.}
La frontiera è costruita a runtime immediatamente prima di ogni
selezione, a partire dai profili energetici correnti. Come coordinata
energetica viene utilizzata la stima LCB descritta nella
Sezione~\ref{sec:esplorazione-ucb}, in modo che le varianti
sotto-campionate beneficino del bonus di esplorazione anche ai fini
del posizionamento sulla frontiera.

L'algoritmo ordina le varianti per energia crescente e percorre la
lista mantenendo il minimo errore osservato: una variante è
Pareto-ottimale se e solo se il suo errore è strettamente inferiore al
minimo corrente. La complessità risultante è $O(n \log n)$, dominata
dall'ordinamento.


% ============================================================
\section{Scalarizzazione e selezione della variante}
\label{sec:scalarizzazione}
% ============================================================

La frontiera di Pareto individua l'insieme dei compromessi ammissibili,
ma non determina quale scegliere: la selezione dipende dall'importanza
relativa delle due dimensioni, codificata dal peso $\lambda$.

\subsection{Normalizzazione e funzione obiettivo}
\label{sec:scalarizzazione:funzione}

Per rendere le due dimensioni confrontabili, energia ed errore sono
normalizzati nell'intervallo $[0,1]$ rispetto al minimo e al massimo
osservati sulla frontiera~$\mathcal{P}$:
%
\begin{align}
    \hat{E}(v) &= \frac{E(v) - E_{\min}}{E_{\max} - E_{\min}}
    \label{eq:norm_energy} \\[4pt]
    \widehat{\mathit{Err}}(v) &=
        \frac{\mathit{Err}(v) - \mathit{Err}_{\min}}
             {\mathit{Err}_{\max} - \mathit{Err}_{\min}}
    \label{eq:norm_error}
\end{align}

La funzione di scalarizzazione lineare è definita come:
%
\begin{equation}
    S(v,\,\lambda) \;=\;
        (1 - \lambda)\,\hat{E}(v) \;+\;
        \lambda\,\widehat{\mathit{Err}}(v)
    \label{eq:scalarizzazione}
\end{equation}

La variante selezionata $v^*$ è quella che minimizza $S$ sulla
frontiera:
%
\begin{equation}
    v^* = \arg\min_{v \in \mathcal{P}}\; S(v,\,\lambda)
    \label{eq:selezione_carbon}
\end{equation}

\medskip
\noindent\textbf{Interpretazione dei casi limite.}
Quando $\lambda = 1$ (rete pulita, CI bassa) la funzione si riduce a
$S = \widehat{\mathit{Err}}(v)$: il sistema seleziona la variante più
accurata. Quando $\lambda = 0$ (rete sporca, CI alta) si ha
$S = \hat{E}(v)$: il sistema seleziona la variante a consumo minimo.
Per valori intermedi la selezione bilancia le due dimensioni in modo
continuo e proporzionale al segnale ambientale.

Per costruzione, $v^*$ appartiene sempre alla frontiera di Pareto ed è
quindi non dominata: nessuna alternativa è contemporaneamente più
efficiente e più accurata.


% ============================================================
\section{Formalizzazione della logica decisionale}
\label{sec:carbon:algoritmo}
% ============================================================

La logica complessiva dello scheduler carbon-aware è formalizzata
nell'Algoritmo~\ref{alg:carbon-select}, che integra le fasi descritte
nelle sezioni precedenti.

\begin{algorithm}[htbp]
\caption{\textsc{SelectParetoVariant}: selezione carbon-aware}
\label{alg:carbon-select}
\begin{algorithmic}[1]
\Function{SelectParetoVariant}{$r$ : \texttt{Request}}
    \State $\textit{base} \gets r.\texttt{Fun}$
    \State $V \gets \Call{GetFunctionsByLogicalName}
                        {\textit{base}.\texttt{LogicalName}}$
    \If{$V = \emptyset$}
        \State \Return $(\textit{base},\;
               \text{``fallback-base''})$
    \EndIf
    \State
    \State \Comment{Valutazione delle varianti (energia LCB + errore)}
    \State $\textit{Eval} \gets [\,]$
    \For{$v \in V$}
        \State $\textit{warm} \gets \Call{HasWarmContainer}{v}$
        \State $E(v) \gets \Call{EnergyLCB}{v,\;\textit{warm},
                   \;\beta,\;\texttt{minSamples}}$
        \State $\mathit{Err}(v) \gets \Call{ErrorScore}{v}$
        \State $\textit{Eval}.\textsc{Append}(v,\;E(v),
                   \;\mathit{Err}(v))$
    \EndFor
    \State
    \State \Comment{Frontiera di Pareto --- $O(n \log n)$}
    \State $\mathcal{P} \gets \Call{ParetoFilter}{\textit{Eval}}$
    \If{$|\mathcal{P}| = 1$}
        \State \Return $(\mathcal{P}[0],\;
               \text{``pareto-singleton''})$
    \EndIf
    \State
    \State \Comment{Calcolo di $\lambda$ dalla CI corrente}
    \State $\mathrm{CI} \gets \Call{GetCarbonIntensity}
                   {r.\texttt{Zone}}$
    \State $(\mu_z,\, s_z,\, k) \gets \Call{GetSigmoidParams}
                   {r.\texttt{Zone}}$
    \State $\lambda \gets \dfrac{1}{1 + \exp\!\bigl(
               k\,(\mathrm{CI} - \mu_z)\,/\,s_z\bigr)}$
    \State
    \State \Comment{Normalizzazione e scalarizzazione}
    \State $E_{\min},\; E_{\max} \gets \Call{MinMax}
               {[\,E(v) : v \in \mathcal{P}\,]}$
    \State $\mathit{Err}_{\min},\; \mathit{Err}_{\max} \gets
           \Call{MinMax}
               {[\,\mathit{Err}(v) : v \in \mathcal{P}\,]}$
    \State $v^* \gets \displaystyle\arg\min_{v \in \mathcal{P}}\;
           \bigl[(1-\lambda)\,\hat{E}(v) +
                  \lambda\,\widehat{\mathit{Err}}(v)\bigr]$
    \State \Return $(v^*,\; \text{``pareto-scalarisation''})$
\EndFunction
\end{algorithmic}
\end{algorithm}

L'algoritmo evidenzia la separazione tra le fasi di raccolta dei
profili, costruzione della frontiera, calcolo di $\lambda$ e
scalarizzazione. La struttura modulare consente di sostituire ciascun
componente indipendentemente: ad esempio, il modello di normalizzazione
della CI potrebbe essere sostituito con un modello non parametrico senza
modificare la logica di selezione.

La stima energetica utilizzata nella fase di valutazione incorpora il
meccanismo LCB descritto nella Sezione~\ref{sec:esplorazione-ucb}: le
varianti sotto-campionate ricevono una stima ottimisticamente ridotta,
che ne favorisce la selezione e l'accumulo di osservazioni reali.
Le chiamate \textsc{GetCarbonIntensity} e \textsc{GetSigmoidParams}
operano su dati in memoria locale (la CI è aggiornata in background
ogni 15~minuti), evitando latenze di rete nel percorso critico della
richiesta.


% ============================================================
\section{Discussione delle scelte progettuali}
\label{sec:carbon:discussione}
% ============================================================

Un approccio alternativo alla scalarizzazione continua avrebbe potuto
adottare una soglia sulla CI: al di sopra di un livello predefinito,
selezionare sempre la variante a consumo minimo; al di sotto, quella a
errore minimo. Tale schema, pur essendo più semplice, è privo di
continuità: variazioni minime della CI intorno alla soglia producono
transizioni brusche nella selezione. La scalarizzazione con $\lambda$
continuo produce invece una transizione fluida, proprietà desiderabile
nei contesti in cui la CI varia lentamente nel tempo, come nei profili
giornalieri delle reti elettriche.

La calibrazione della sigmoide sulla distribuzione storica locale,
anziché su una scala assoluta, riflette il fatto che il giudizio su una
condizione di CI \emph{alta} o \emph{bassa} è intrinsecamente
contestuale. Ciò introduce una dipendenza dalla qualità dei dati
storici: zone con distribuzioni non rappresentative possono determinare
calibrazioni sub-ottimali. L'impatto di questa limitazione è valutato
nel Capitolo~\ref{cap:RisultatiSperimentali}.

Infine, la funzione di scalarizzazione~\eqref{eq:scalarizzazione} è
generalizzabile a forme non lineari --- quali la scalarizzazione di
Chebyshev o la media $p$-pesata --- senza modifiche all'architettura,
poiché la normalizzazione sulla frontiera Pareto-ottimale garantisce in
ogni caso la non-dominanza della soluzione selezionata.